% =====================================================================
% PROPOSAL --- response-letter text for the price-reporting / closure
% clarification. NOT \input anywhere; the manuscript compiles without it.
%
% Adoption step:
%   1. Decide the home. Variant A (R2.15) if the units point came from
%      referee 2; Variant B (E2) if it came from the editor's letter or a
%      third report. Everything except the heading is identical.
%   2. Paste the chosen variant at the END of revision/response/reviewer2.tex
%      (respectively revision/response/editor.tex).
%   3. Insert the referee's own words into \refcom{...}.
%   4. Check the two numbers against revision/tables/pricing_prices.csv
%      before submission (they are read from it, not from memory: 60
%      price-moving GAMMA/DELTA cells, 10 sign reversals, worst case
%      DELTA-F3-k05-spis at +1.92 % in wage units and -0.01 % in
%      consumption units).
%
% ADOPTED 2026-10-23: wired into response/reviewer2.tex as R2.15 (heading
% 'Units of the reported price changes'). This file stays as the record of the
% wording; only the \refcom quote still has to be filled in.
%
% Written 2026-10-23 (AG) under ADR-0037 (consumption-unit reporting) and
% ADR-0038 (employment-adjustment closure). Cross-checked against
% Corollary cor:peg and Remark rem:empadj in Appendix B.
% =====================================================================

% ---------- Variant A: as a referee point ----------------------------
\subsection*{R2.15 --- Units of the reported price changes}
\refcom{[insert the referee's wording on the normalisation of reported
price changes]}
\response{We thank the referee: this was wrong, and chasing it turned up a
second, more substantive error of ours. The two are related but distinct, and
we have treated them as such.

\emph{Units.} The pricing tables compared nominal magnitudes expressed in two
different units. The mobile closures carry the consumption price index as
numeraire, so their rows were already in consumption units; the wage-pinned
rows fixed the money wage instead, leaving prices in wage units. All pricing
rows are now reported in consumption units: prices are divided by the cell's
own CPI, the GDP deflator is rescaled with them, and the maximum sectoral
deviation is computed on the normalised price vector. Where the CPI sits at
one --- which is the whole demand-only matrix, and therefore every headline
result of the paper --- nothing changes. Where the pricing doors move prices,
it does: of the sixty affected cells ten reverse the sign of the reported
deflator, and in thirty-two of them the correction is larger than the effect
being reported. The extreme case is the distributed-profit, Leontief,
externally financed cell at $\kappa = 0.5$, which read $+1.92$ per cent in
wage units and reads $-0.01$ per cent in consumption units: the entire
"effect" was the change of unit.

\emph{Closure.} The reason the CPI moved at all is that the closure we had
called a ``fixed real wage'' closure does not peg the real wage. It fixes the
money wage as the numeraire, sets the net external transfer to zero and drops
the labour-supply equation, so that employment adjusts. With one production
factor and constant returns, the real wage is not an instrument at all: it is
fixed by technology, $w/\mathcal{P}(p) = 1/\mathcal{P}(\pi(A))$, so a peg on it
is redundant at the calibration point and infeasible away from it. We have
renamed the closure \emph{employment adjustment}, stated its equations exactly
as they are solved, and added the identity with its proof (Appendix~B,
the no-real-wage-peg corollary and the remark on what the closure implements). No estimate in the paper
changes: the kernel never solved the peg, and real outcomes are invariant to
the nominal scale. What changes is the label, the interpretation --- in
particular we no longer claim that workers' purchasing power is held fixed
there --- and the units.

The two corrections cannot substitute for one another, which is why we report
both: renormalising cannot restore a constant real wage, and relabelling
cannot repair mixed units. Reporting the wage-pinned rows in consumption units
is in fact what makes the movement of the real wage visible rather than hidden
inside the reporting convention.}
\loc{Sec.~4 (numeraire and the reporting rule); Secs.~5--6 (the affected
tables); Appendix~A (the nesting claims, now scoped to the demand-only
slice); Appendix~B (col.\ the no-real-wage-peg corollary and the accompanying remark).}

% ---------- Variant B: as a self-disclosed correction -----------------
% \subsection*{E2 --- Corrections we made on our own initiative}
% \refcom{[the editor's letter does not raise this point]}
% \response{[same body as Variant A, with ``We thank the referee'' replaced by
%   ``Beyond the points raised by the referees, we report two corrections of
%   our own. They concern how the fixed-wage closure is described and in which
%   units the reported price changes are expressed. Neither changes a real
%   result; one changes the interpretation of a whole column of our matrix.'']}
% \loc{[as Variant A]}
